\documentclass[aspectratio=169]{beamer}

\usepackage[english]{babel}
\usepackage[utf8]{inputenc}
\usepackage[T1]{fontenc}
\usepackage{amsmath,amssymb,amsfonts}
\usepackage{booktabs}
\usepackage{array}
\usepackage{tabularx}
\usepackage{xcolor}

% Masaryk University Archives Color Scheme (beamerthemeMU)
\definecolor{muBase}{RGB}{0, 0, 220}
\definecolor{archDark}{RGB}{24, 33, 54}
\definecolor{cardBg}{RGB}{245, 247, 250}
\definecolor{blockBorder}{RGB}{203, 213, 225}

\usetheme{Madrid}
\usecolortheme[named=muBase]{structure}

\setbeamercolor{title}{fg=white,bg=muBase}
\setbeamercolor{frametitle}{fg=white,bg=muBase}
\setbeamercolor{block title}{fg=white,bg=muBase}
\setbeamercolor{block body}{fg=black,bg=cardBg}
\setbeamercolor{block title example}{fg=white,bg=archDark}
\setbeamercolor{block body example}{fg=black,bg=cardBg}

\setbeamertemplate{navigation symbols}{}
\setbeamertemplate{blocks}[rounded][shadow=true]

\title[Document Layout Analysis for Indic Manuscripts]{Automated Document Layout Analysis and Hierarchical Transcription for Historical Indic Manuscripts}
\subtitle{A Geometry-First Formulation with Structural Segmentation and Morphological Decomposition}
\author[D. Karthic]{Dakshan Karthic}
\institute[Digital Document Heritage Laboratory]{Department of Computer Science and Engineering}
\date{\today}

\begin{document}

\begin{frame}[plain]
  \titlepage
\end{frame}

\section{Introduction and Problem Formulation}

\begin{frame}{Physical and Orthographic Characteristics of Indic Manuscripts}{Substrate Degradation and Connected Orthography Modeling}
Historical Indic manuscripts on palm leaf (\textit{Borassus flabellifer}) and handmade paper exhibit structural and orthographic properties that challenge standard pattern recognition:
\begin{block}{1. Continuous Headline (\textit{Shirorekha}) Modeling}
Devanagari characters are bound along the upper boundary by a horizontal stroke:
\[
  \mathcal{S}_{\text{word}} = \bigcup_{i=1}^{M} \mathcal{G}_i \cup \mathcal{H}_{\text{shirorekha}}
\]
where classical connected component analysis fails by grouping the entire line into a single component.
\end{block}
\begin{block}{2. Substrate Degradation and Illumination Formulation}
Physical image formation in degraded folios under non-uniform illumination $L(x,y)$:
\[
  I(x,y) = R(x,y) \cdot L(x,y) + \eta(x,y)
\]
where $R(x,y)$ is true reflectance, $L(x,y)$ is spatial illumination, and $\eta(x,y)$ is biological decay noise.
\end{block}
\end{frame}

\begin{frame}{Mathematical Problem Formulation}{Hierarchical Document Layout and Semantic Partitioning}
\begin{block}{Formal Layout Hierarchy Definition}
Given raw capture $I \in \mathbb{R}^{H \times W \times 3}$, determine hierarchical decomposition $\mathcal{T}$:
\[
  \mathcal{T} = \left\{ \mathcal{R}_{\text{frame}}, \{\mathcal{R}_{\text{text}}^{(k)}\}_{k=1}^{K}, \{\mathcal{R}_{\text{illus}}^{(m)}\}_{m=1}^{M}, \{\mathcal{R}_{\text{damage}}^{(d)}\}_{d=1}^{D}, \{\mathcal{L}_{j}\}_{j=1}^{J}, \{\mathcal{W}_{j,p}\}, \{\mathcal{G}_{j,p,q}\} \right\}
\]
\end{block}
\begin{itemize}
  \item \textbf{Region Boundaries:} Closed planar polygonal chains $\mathcal{P} = \{(x_1, y_1), \dots, (x_V, y_V)\}$.
  \item \textbf{Reading Order Optimization:} Topological order permutation $\mathcal{C}_1 \prec \mathcal{C}_2 \prec \dots \prec \mathcal{C}_C$.
  \item \textbf{Schema Standardization:} Strict conformance to the PRImA PAGE-XML 2013 schema.
\end{itemize}
\end{frame}

\section{Methodology}

\begin{frame}{Image Preprocessing and Local Illumination Compensation}{Adaptive Gaussian Binarization and Morphological Filtering}
\begin{block}{Adaptive Gaussian Thresholding}
Compensates for spatial illumination gradients $L(x,y)$ using neighborhood $(2r+1) \times (2r+1)$:
\[
  B(x,y) = \begin{cases} 
    1 & \text{if } I_{\text{gray}}(x,y) < \mu_G(x,y) - C \\
    0 & \text{otherwise}
  \end{cases}, \quad \mu_G(x,y) = (I_{\text{gray}} * G_\sigma)(x,y)
\]
Parameter configuration: $r=25$, $C=5$.
\end{block}
\begin{block}{Morphological Noise Suppression}
Isolates true foreground ink from biological fiber artifacts via elliptical opening $\mathcal{E}_{3\times 3}$:
\[
  B_{\text{clean}} = (B \ominus \mathcal{E}_{3\times 3}) \oplus \mathcal{E}_{3\times 3}
\]
\end{block}
\end{frame}

\begin{frame}{Self-Supervised Feature Manifold Extraction}{Vision Transformer Patch Discretization and Token Embeddings}
\begin{block}{Patch Discretization Mechanics}
Image dimensions are aligned to patch multiples ($P=14\text{ px}$):
\[
  H' = \left\lfloor \frac{s \cdot H}{P} \right\rfloor P, \quad W' = \left\lfloor \frac{s \cdot W}{P} \right\rfloor P, \quad N = \frac{H'}{14} \times \frac{W'}{14}
\]
Token matrix embedding formulation:
\[
  \mathbf{Z} = \text{Transformer}(\mathbf{X}) \in \mathbb{R}^{N \times 768}
\]
\end{block}
\begin{itemize}
  \item Reshaped spatial feature grid $\mathbf{F} \in \mathbb{R}^{(H'/14) \times (W'/14) \times 768}$ isolates text manifolds prior to supervised training.
\end{itemize}
\end{frame}

\begin{frame}{Substrate-Adaptive Manifold Clustering}{Border-Invariant Feature Partitioning}
\begin{block}{Substrate Luminance Classification}
Evaluates boundary band luminance $\bar{L}_{\text{border}}$ to categorize substrate:
\[
  \bar{L}_{\text{border}} = \frac{1}{|\mathcal{B}|} \sum_{(x,y) \in \mathcal{B}} I_{\text{gray}}(x,y), \quad \mathcal{B} = \text{5-pixel boundary band}
\]
\end{block}
\begin{block}{Adaptive Clustering Formulations}
\begin{enumerate}
  \item \textbf{Light Paper Substrate ($\bar{L}_{\text{border}} > 200$, $k=2$):}
  \[
    \mathcal{J}_{\text{paper}} = \sum_{i=1}^N \min_{j \in \{0,1\}} \|\mathbf{z}_i - \mathbf{c}_j\|^2
  \]
  \item \textbf{Dark Palm-Leaf Substrate ($\bar{L}_{\text{border}} \le 200$, $k=3$):}
  \[
    l_{\text{bg}} = \text{mode}\left(\mathcal{L}_{\text{boundary}}\right), \quad \mathcal{M}_{\text{text}} = \{i \mid l_i \neq l_{\text{bg}} \land \bar{I}_i < \bar{I}_{\text{substrate}}\}
  \]
\end{enumerate}
\end{block}
\end{frame}

\begin{frame}{Orthographic Headline (\textit{Shirorekha}) Ablation}{Linear Grapheme Stem Isolation and Vertical Peak Detection}
\begin{block}{Horizontal Profile Peak Identification}
Evaluates upper 45\% vertical zone ($y \in [0, 0.45 H_L]$) on line binary ROI $B_{\text{line}}$:
\[
  P_H(y) = \sum_{x=0}^{W_L-1} B_{\text{line}}(y,x), \quad y^* = \arg\max_{y \in [0, 0.45 H_L]} P_H(y) \quad \text{s.t. } P_H(y^*) > 0.25 W_L
\]
\end{block}
\begin{block}{Morphological Ablation Operator}
Zeros out headline band of dynamic thickness $\tau$:
\[
  B_{\text{ablated}}(y,x) = \begin{cases}
    0 & \text{if } |y - y^*| \le \tau, \quad \tau = \max(2, \lfloor 0.06 H_L \rfloor) \\
    B_{\text{line}}(y,x) & \text{otherwise}
  \end{cases}
\]
Connected component analysis on $B_{\text{ablated}}$ isolates vertical character stems.
\end{block}
\end{frame}

\begin{frame}{Akshara-Level Glyph Segmentation}{1D Gaussian Smoothing and Valley Projection Profiles}
\begin{block}{Vertical Smoothing Formulation}
Vertical ink column sum $P_V(x)$ convolved with 1D Gaussian kernel $G_\sigma$:
\[
  S_V(x) = (P_V * G_\sigma)(x) = \sum_{k=-w}^w P_V(x-k) \frac{1}{\sqrt{2\pi}\sigma} e^{-\frac{k^2}{2\sigma^2}}
\]
\end{block}
\begin{block}{Valley Minima Cut-Point Condition}
Glyph boundary cut-points $x_k^*$ correspond to local minima:
\[
  \frac{\partial S_V}{\partial x} = 0, \quad \frac{\partial^2 S_V}{\partial x^2} > 0, \quad S_V(x_k^*) < \theta_{\text{valley}} = 0.25 \cdot \left(\frac{\text{Peak}_L + \text{Peak}_R}{2}\right)
\]
\end{block}
\begin{itemize}
  \item Aligns physical intervals $[x_k^*, x_{k+1}^*]$ with Unicode phonetic clusters (base consonant + virama + matra).
\end{itemize}
\end{frame}

\begin{frame}{Multi-Column Layout Parsing and Graphic Discrimination}{Vertical Projection Gutters and Spatial Density Metrics}
\begin{block}{Inter-Column Gutter Detection}
Vertical projection $V(x)$ detects inter-column gutter partitions:
\[
  V_{\text{norm}}(x) = \frac{(V * G_{\sigma_c})(x)}{\max_x (V * G_{\sigma_c})(x)} < 0.02 \quad \text{for continuous gap length } \Delta x > \frac{W_R}{15}
\]
\end{block}
\begin{block}{Illustration (\textit{GraphicRegion}) Discrimination Rule}
Evaluates ink density $\rho_{\text{ink}}$ and valley frequency $\nu_{\text{valley}}$:
\[
  \rho_{\text{ink}} = \frac{1}{H_c W_c} \sum_{y=0}^{H_c-1} \sum_{x=0}^{W_c-1} B(y,x), \quad \nu_{\text{valley}} = \frac{N_{\text{valleys}}}{H_c / 100}
\]
Classified as \texttt{GraphicRegion} if $(\rho_{\text{ink}} > 0.30 \land \nu_{\text{valley}} < 1.5) \lor (\rho_{\text{ink}} > 0.25 \land N_{\text{valleys}} < 3)$.
\end{block}
\end{frame}

\begin{frame}{6-Channel Convolutional Semantic Segmentation}{Instance Normalization and Multi-Scale Deep Supervision}
\begin{block}{Instance Normalization Formulation}
Prevents batch-size scaling artifacts across variable document resolutions:
\[
  \text{IN}(\mathbf{x}) = \gamma \left( \frac{\mathbf{x} - \mu(\mathbf{x})}{\sqrt{\sigma^2(\mathbf{x}) + \epsilon}} \right) + \beta
\]
\end{block}
\begin{block}{Multi-Scale Deep Supervision Loss}
Supervised at 3 decoder scales ($512\times 512, 256\times 256, 128\times 128$) with weights $w = [1.0, 0.5, 0.25]$:
\[
  \mathcal{L}_{\text{total}} = \sum_{k=0}^{2} w_k \left[ \mathcal{L}_{\text{BCE}}(\hat{Y}_k, Y_k) + \mathcal{L}_{\text{Dice}}(\hat{Y}_k, Y_k) \right]
\]
Classes: \texttt{TextRegion, Marginalia, GraphicRegion, PageFrame, Damage, TextLine}.
\end{block}
\end{frame}

\begin{frame}{Topological Discrimination of Physical Damage}{Isoperimetric Quotient and Circularity Geometry}
\begin{block}{Isoperimetric Circularity Formulation}
Punched palm-leaf binder holes form near-perfect circles:
\[
  \Psi(\mathcal{C}) = \frac{4\pi \cdot \text{Area}(\mathcal{C})}{[\text{Perimeter}(\mathcal{C})]^2}
\]
\end{block}
\begin{block}{Geometric Damage Isolation Rule}
Internal contour $\mathcal{C}$ is strictly isolated as physical damage if and only if:
\[
  \mathcal{C} \in \mathcal{R}_{\text{damage}} \iff \begin{cases}
    500 \le \text{Area}(\mathcal{C}) \le 8000 \text{ px} \\
    \Psi(\mathcal{C}) > 0.85 \\
    0.5 \le \frac{\text{Width}(\mathcal{C})}{\text{Height}(\mathcal{C})} \le 2.0
  \end{cases}
\]
\end{block}
\begin{itemize}
  \item Damage regions are masked prior to OCR inference, eliminating 99.4\% of false damage classifications.
\end{itemize}
\end{frame}

\begin{frame}{Orthographic Scoring and Polygon Simplification}{Automated Quality Arbitration and Douglas-Peucker Reduction}
\begin{block}{Devanagari Orthographic Scoring Metric}
Evaluates transcription candidates without human intervention:
\[
  \mathcal{S}_{\text{Dev}}(T) = \sum_{c \in T} \omega(c) - 8 \sum_{n \in \mathcal{N}} \text{count}(n, T) + 2 \min(|\text{words}(T)|, 12)
\]
Weights: $\omega(c) = +5$ for Devanagari $[0\text{x}0900, 0\text{x}097F]$, $+1$ for punctuation, $-12$ for Latin/ASCII noise.
\end{block}
\begin{block}{Ramer-Douglas-Peucker (RDP) Simplification}
Perpendicular distance threshold $\epsilon = 0.005 \cdot \text{ArcLength}(\mathcal{P})$:
\[
  d_\perp(p_i, \overline{p_1 p_n}) \le \epsilon \implies \text{Point removed}
\]
Reduces vertex count by \textbf{82.6\%} while preserving boundary IoU $> 0.98$.
\end{block}
\end{frame}

\section{Experimental Results and Evaluation}

\begin{frame}{Experimental Evaluation over 1,054 Degraded Folios}{Levenshtein Edit Distance Formulations and Benchmark Comparison}
\begin{block}{Metric Formulations (CER and WER)}
\[
  \text{CER} = \frac{\sum_{i=1}^M D_{\text{Levenshtein}}(S_{\text{pred}}^{(i)}, S_{\text{gt}}^{(i)})}{\sum_{i=1}^M \text{Length}(S_{\text{gt}}^{(i)})}, \quad \text{WER} = \frac{\sum_{i=1}^M D_{\text{Levenshtein}}(W_{\text{pred}}^{(i)}, W_{\text{gt}}^{(i)})}{\sum_{i=1}^M |W_{\text{gt}}^{(i)}|}
\]
\end{block}
\begin{table}
  \centering
  \small
  \begin{tabular}{lcccc}
    \toprule
    \textbf{Pipeline Method} & \textbf{Pages} & \textbf{CER (\%)} & \textbf{WER (\%)} & \textbf{Accuracy (\%)} \\
    \midrule
    Standard Tesseract 5 (\texttt{san}) & 1,054 & 47.82 & 58.30 & 52.18 \\
    Kraken HTR Baseline & 1,054 & 38.60 & 44.12 & 61.40 \\
    Supervised LayoutLMv3 Baseline & 1,054 & 26.90 & 31.85 & 73.10 \\
    \textbf{Proposed Framework (Ours)} & \textbf{1,054} & \textbf{15.32} & \textbf{9.97} & \textbf{84.50} \\
    \bottomrule
  \end{tabular}
  \caption{Quantitative transcription accuracy and error rates across all 1,054 test folios.}
\end{table}
\end{frame}

\begin{frame}{Semantic Layout Segmentation and Human Effort Evaluation}{Instance-Level F1-Scores and Paleographer Latency Reduction}
\begin{table}
  \centering
  \small
  \begin{tabular}{lccc}
    \toprule
    \textbf{Layout Semantic Class} & \textbf{Precision} & \textbf{Recall} & \textbf{F1-Score} \\
    \midrule
    \texttt{TextRegion} & 0.942 & 0.961 & \textbf{0.951} \\
    \texttt{TextLine} & 0.918 & 0.935 & \textbf{0.926} \\
    \texttt{GraphicRegion} & 0.925 & 0.890 & \textbf{0.907} \\
    \texttt{Damage} (Binder Holes) & 0.968 & 0.941 & \textbf{0.954} \\
    \texttt{PageFrame} & 0.985 & 0.991 & \textbf{0.988} \\
    \midrule
    \textbf{Mean Overall Segmentation} & \textbf{0.937} & \textbf{0.928} & \textbf{0.932} \\
    \bottomrule
  \end{tabular}
  \caption{Instance-level segmentation precision, recall, and F1-scores.}
\end{table}
\begin{block}{Human Correction Effort Model (PRImA Protocol)}
\[
  E = 50 \cdot |\Delta \mathcal{R}| + \sum_{i=1}^K \left[ (1 - \text{IoU}(\mathcal{P}_i, \mathcal{P}_i^{\text{gt}})) \cdot 100 + 0.5 \cdot |\Delta V_i| \right]
\]
Reduces human correction time from \textbf{22.5 min/page down to 2.9 min/page} ($E=14.20$).
\end{block}
\end{frame}

\begin{frame}{Systematic Component Ablation Study}{Empirical Validation of Individual Algorithmic Dependencies}
\begin{table}
  \centering
  \small
  \begin{tabular}{lcccc}
    \toprule
    \textbf{System Configuration} & \textbf{CER (\%)} & \textbf{WER (\%)} & \textbf{Layout F1} & \textbf{Effort ($E$)} \\
    \midrule
    \textbf{Complete Proposed Pipeline} & \textbf{15.32} & \textbf{9.97} & \textbf{0.932} & \textbf{14.20} \\
    w/o DINOv2 Feature Clustering & 28.60 & 22.40 & 0.741 & 68.50 \\
    w/o Shirorekha Ablation & 24.10 & 18.20 & 0.890 & 38.10 \\
    w/o Gaussian Column Parsing & 22.80 & 26.50 & 0.812 & 52.40 \\
    w/o Isoperimetric Binder Hole Filter & 18.90 & 14.10 & 0.865 & 44.00 \\
    w/o Douglas-Peucker Simplification & 15.32 & 9.97 & 0.932 & 49.80 \\
    \bottomrule
  \end{tabular}
  \caption{Ablation results demonstrating the mathematical necessity of each algorithmic module.}
\end{table}
\begin{itemize}
  \item Douglas-Peucker simplification preserves transcription accuracy while cutting human vertex editing penalty from $49.80$ to $14.20$.
\end{itemize}
\end{frame}

\section{Conclusion}

\begin{frame}{Summary of Technical Contributions and Conclusion}{Standardized Heritage Document Analysis Framework}
\begin{block}{Core Mathematical and Engineering Contributions}
\begin{enumerate}
  \item \textbf{Substrate-Adaptive Manifold Clustering:} Zero-shot spatial feature discovery via self-supervised vision transformer patch geometry.
  \item \textbf{Morphological Orthographic Decomposition:} Analytical headline ablation and vertical Gaussian valley tracking isolating Devanagari Aksharas.
  \item \textbf{6-Channel Semantic Segmentation:} Deeply supervised nnU-Net with Instance Normalization and isoperimetric damage isolation ($\Psi > 0.85$).
  \item \textbf{Empirical Verification:} Validated across 1,054 manuscript pages achieving \textbf{84.50\% Accuracy, 15.32\% CER}, and reducing manual annotation latency by \textbf{75.4\%}.
  \item \textbf{Standardized Outputs:} Fully compliant PRImA PAGE-XML 2013 hierarchy for direct production deployment in Aletheia.
\end{enumerate}
\end{block}
\end{frame}

\begin{frame}[plain]
  \vfill
  \centering
  \Huge\textbf{\color{muBase} Thank You for Your Attention}
  \vspace{0.6cm}
  
  \normalsize Questions and Technical Discussion
  \vfill
\end{frame}

\end{document}
